%% ===== Methods (no word limit; full math here) =====
%% Nature/NBT: Methods is a dedicated section at the END of the main article (publishes
%% online with the paper), NOT in the Supplementary Information. Methods references
%% continue the numbering after the main reference list. Subsections are unnumbered
%% (\subsection*) per Nature style. Data/Code availability + Reporting summary live here.
\section{Methods\cb{no limit}}\label{sec:methods}

\subsection*{Ontology and knowledge-graph construction}
\textit{[FILLER.]} Experiment ontology (genotype, phenotype, environment, experiment
metadata); annotation of literature and public datasets; ingestion into the Neo4j
knowledge graph; conversion, deduplication, and aggregation into HeteroData records
(Suppl.\ Figs.\ \ref{fig:data-model}, \ref{fig:conversion-dedup-agg}).

\subsection*{Datasets}
\textit{[FILLER.]} Source datasets and versions (e.g.\ Costanzo, Kuzmin trigenic;
knockout expression; single-knockout morphology); inclusion/exclusion criteria; dataset
sizes and label provenance.

\subsection*{Cell representation and the perturbation operator}
\textit{[FILLER.]} Shared wildtype reference cell; each strain as a perturbation operator
over the reference; multimodal node/edge features (genome sequence, gene-interaction
networks, metabolism, environment); DNA sequence selection and tokenization
(Suppl.\ Figs.\ \ref{fig:dna-selection}, \ref{fig:dna-tokenization}).

\subsection*{Data splits and preprocessing}
\textit{[FILLER.]} Train/validation/test partitioning (and any gene- or
interaction-disjoint splits); label normalization; handling of missing labels.

\subsection*{Cell Graph Transformer architecture}
\textit{[FILLER.]} Graph-constrained multi-head attention; the equivariant perturbation
operator mapping reference embedding to perturbed state; multitask decoder heads
(Fig.~\ref{fig:torchcell}c). Layer counts, hidden dimensions, and parameter budget.

\subsection*{Graph-regularized attention and training objective}
\textit{[FILLER.]} Attention-prior alignment loss aligning attention to network priors;
multitask loss weighting; full mathematical formulation.

\subsection*{Model training}
\textit{[FILLER.]} Optimizer, learning-rate schedule, batch size, early stopping;
hardware (GPUs, distributed/DDP), wall-clock; software versions; random seeds.

\subsection*{Baselines}
\textit{[FILLER.]} DANGO, DCell, and Yeast9 flux balance analysis; classical-ML baselines
(Suppl.\ Fig.\ S1); matched training data and evaluation protocol.

\subsection*{Evaluation metrics}
\textit{[FILLER.]} Pearson correlation ($r$) and any secondary metrics; aggregation
across folds; definition of error bars (see Statistics).

\subsection*{Gene--gene interaction prediction}
\textit{[FILLER.]} Trigenic interaction prediction setup; performance vs.\ interaction
magnitude; ablations (Fig.~\ref{fig:ggi}).

\subsection*{Expression and morphology prediction}
\textit{[FILLER.]} Knockout transcriptome and single-knockout morphology tasks; shared
latent embedding; cross-phenotype transfer (Fig.~\ref{fig:multitask}).

\subsection*{CGT-guided strain design}
\textit{[FILLER.]} Objective and candidate enumeration for beta-carotene and betaxanthin
production; ranking and selection of recommended gene deletions (Fig.~\ref{fig:metabolism}).

\subsection*{Strain construction and cultivation}
\textit{[FILLER.]} Strains, plasmids, and genome edits; the TorchCell + UIUC iBioFoundry
design--build--test--learn loop; cultivation and fermentation conditions.

\subsection*{Metabolite quantification}
\textit{[FILLER.]} Sampling, extraction, and analytical method (HPLC/LC--MS); calibration
and titer quantification.

\subsection*{Inference at scale}
\textit{[FILLER.]} Large-scale inference pipeline and resources (Suppl.\ Fig.\ S6).

\subsection*{Statistics}
\textit{[FILLER.]} Statistical tests used and whether one- or two-sided; sample sizes
($n$); definition of center and dispersion; box-plot conventions (median, IQR hinges,
whiskers $\le 1.5\times$IQR); multiple-comparison handling; significance thresholds.

\subsection*{Reporting summary}
Further information on research design is available in the Nature Portfolio Reporting
Summary linked to this article.

\subsection*{Data availability}
\textit{[FILLER.]} Datasets are available on Hugging Face
(\url{https://huggingface.co/...}); the Neo4j knowledge graph is hosted at NCSA
(\url{...}); trained model weights are available at \url{...}. Source data are provided
with this paper.

\subsection*{Code availability}
\textit{[FILLER.]} TorchCell and the Cell Graph Transformer are open source at
\url{https://github.com/Mjvolk3/torchcell} (license: ...), with a release archived at
Zenodo (\url{https://doi.org/...}).
